\documentclass[12pt,a4paper]{article}
\usepackage[T1]{fontenc}
\usepackage[utf8]{inputenc}
\usepackage[spanish]{babel}
\usepackage{amsmath, amssymb}
\usepackage{graphicx}
\usepackage{hyperref}
\usepackage{listings}
\usepackage{xcolor}
\usepackage{array}
\usepackage{booktabs}
\usepackage[margin=2.5cm]{geometry}

% Configuración de listings para código Python
\lstset{
    language=Python,
    basicstyle=\ttfamily\small,
    keywordstyle=\color{blue},
    commentstyle=\color{gray},
    stringstyle=\color{red},
    breaklines=true,
    frame=single,
    rulecolor=\color{lightgray},
    showstringspaces=false,
    tabsize=4,
}

\title{\textbf{Referencia Técnica: Librería \texttt{fatigueset-lib}}}
\author{Pipeline de Procesamiento del Dataset FatigueSet}
\date{\today}

\begin{document}

\maketitle

\tableofcontents
\newpage

% ========================================================================
\section{Introducción}
% ========================================================================

La librería \texttt{fatigueset-lib} es un módulo Python especializado en la carga, validación, procesamiento y extracción de características (\textit{features}) del dataset \textbf{FatigueSet}. Este dataset combina múltiples sensores fisiológicos con tareas cognitivas para analizar fatiga física y mental en 12 participantes.

\subsection{Propósito}

\texttt{fatigueset-lib} implementa un pipeline completo y modular que:
\begin{enumerate}
    \item \textbf{Carga} datos desde la estructura de carpetas del dataset FatigueSet
    \item \textbf{Valida} la integridad de los datos (valores nulos, duplicados, tipos)
    \item \textbf{Procesa} series temporales multimodales (sincronización, segmentación por fases)
    \item \textbf{Extrae} características estadísticas y especializadas
    \item \textbf{Agrega} los datos en un dataset listo para Machine Learning (108 filas)
\end{enumerate}

\subsection{Alcance del Dataset}

\begin{itemize}
    \item \textbf{Participantes}: 12 (IDs: 01--12)
    \item \textbf{Sesiones por participante}: 3 (01, 02, 03)
    \item \textbf{Fases de medición}: 3 por sesión
        \begin{itemize}
            \item M1: Baseline (\textit{línea base})
            \item M2: Post-ejercicio físico
            \item M3: Post-carga cognitiva
        \end{itemize}
    \item \textbf{Dataset agregado final}: $12 \times 3 \times 3 = 108$ filas
    \item \textbf{Sensores}: 6 modalidades (pecho, muñeca, frente EEG, oído, etc.)
    \item \textbf{Tareas cognitivas}: N-Back y Choice Reaction Time (CRT)
\end{itemize}

\subsection{Relación con Series Temporales}

Esta librería implementa conceptos clave del documento \textit{``Procesamiento de Datos de Series Temporales''} (véase \texttt{../series-temporales/}), incluyendo:
\begin{itemize}
    \item \textbf{Sincronización} de sensores con distintas frecuencias de muestreo
    \item \textbf{Segmentación temporal} en ventanas (fases M1, M2, M3)
    \item \textbf{Normalización} correcta sin data leakage
    \item \textbf{Extracción de features} estadísticas por ventana
\end{itemize}

% ========================================================================
\section{Descripción General de la Arquitectura}
% ========================================================================

La librería está organizada en 5 módulos principales, cada uno responsable de una etapa del pipeline:

\begin{center}
\begin{tabular}{|l|l|p{6cm}|}
\toprule
\textbf{Módulo} & \textbf{Clase Principal} & \textbf{Responsabilidad} \\
\midrule
\texttt{loader.py} & \texttt{FatigueSetLoader} & Carga de datos desde el dataset \\
\texttt{processor.py} & \texttt{FatigueSetProcessor} & Procesamiento de series temporales \\
\texttt{features.py} & \texttt{FeatureExtractor} & Extracción de características \\
\texttt{validators.py} & \texttt{FatigueSetValidator} & Validación de integridad \\
\texttt{utils.py} & (Funciones) & Funciones auxiliares \\
\bottomrule
\end{tabular}
\end{center}

\subsection{Flujo de datos típico}

\begin{enumerate}
    \item \textbf{Carga}: \texttt{FatigueSetLoader.cargar\_todo()} $\rightarrow$ Dict con todos los DataFrames
    \item \textbf{Validación}: \texttt{FatigueSetValidator.validar\_todos()} $\rightarrow$ Reporte de integridad
    \item \textbf{Procesamiento}: \texttt{FatigueSetProcessor.construir\_dataset\_ml()} $\rightarrow$ Dataset de 108 filas
    \item \textbf{Features}: \texttt{FeatureExtractor.features\_*} $\rightarrow$ Estadísticas especializadas por modalidad
\end{enumerate}

% ========================================================================
\section{Módulo \texttt{loader.py}: Carga de Datos}
% ========================================================================

\subsection{Propósito}

El módulo \texttt{loader.py} proporciona funciones y clases para leer archivos CSV del dataset FatigueSet y construir DataFrames manipulables.

\subsection{Constantes globales}

\begin{lstlisting}
PARTICIPANTES = ['01', '02', ..., '12']
SESIONES = ['01', '02', '03']
FASES_MEDICION = {
    0: 'M1_baseline',
    1: 'M2_post_ejercicio',
    2: 'M3_post_fatiga_mental'
}
INTENSIDAD_NUM = {'low': 1, 'medium': 2, 'high': 3}
CANALES_EEG = ['TP9', 'AF7', 'AF8', 'TP10']
\end{lstlisting}

\subsection{Clase \texttt{DataLoader}}

Cargador genérico de archivos CSV.

\subsubsection{Método \texttt{load\_csv(path)}}

Lee un archivo CSV y devuelve un \texttt{pd.DataFrame}.

\begin{lstlisting}
df = loader.load_csv('datos/archivo.csv')
\end{lstlisting}

\subsubsection{Método \texttt{load\_all(dir\_path)}}

Lee todos los archivos \texttt{.csv} de un directorio como lista de DataFrames.

\subsection{Clase \texttt{FatigueSetLoader}}

Cargador especializado para datos de FatigueSet.

\subsubsection{Inicialización}

\begin{lstlisting}
loader = FatigueSetLoader(
    base_path='fatigueset',  # Ruta del dataset
    participantes=['01', '02', '03'],  # Subconjunto opcional
    sesiones=['01', '02']  # Subconjunto opcional
)
\end{lstlisting}

\subsubsection{Métodos de Metadata e Intensidades}

\begin{itemize}
    \item \texttt{cargar\_metadata()}: Lee \texttt{fatigueset/metadata.csv}. Mapeo de participantes a intensidades (low/medium/high).
    \item \texttt{construir\_intensidad\_map()}: Genera diccionario \texttt{\{pid: \{sesion: intensidad\}\}}.
    \item \texttt{get\_intensidad(participante, sesion)}: Retorna intensidad de una sesión.
\end{itemize}

\subsubsection{Métodos de Carga Individual}

\begin{itemize}
    \item \texttt{cargar\_csv(participante, sesion, archivo, add\_meta=True)}: Carga un archivo específico con metadatos opcionales.
    \item \texttt{cargar\_todos(archivo, verbose=True)}: Carga un archivo de todos los participantes y sesiones configurados.
\end{itemize}

\subsubsection{Métodos especializados por sensor}

\begin{lstlisting}
df_fatiga  = loader.cargar_fatiga()       # exp_fatigue.csv
df_chest   = loader.cargar_chest()        # chest_physiology_summary.csv
df_markers = loader.cargar_markers()      # exp_markers.csv

df_wrist = loader.cargar_wrist()  # Dict con hr, eda, bvp, temp, acc, ibi
df_eeg   = loader.cargar_eeg()    # Dict con alpha, beta, delta, gamma, theta
df_ear   = loader.cargar_ear_ppg()  # Dict con left, right

df_nback = loader.cargar_nback()        # exp_nback.csv
df_crt   = loader.cargar_crt()          # exp_crt.csv
\end{lstlisting}

\subsubsection{Método \texttt{cargar\_todo()}}

Carga \textbf{todas las señales disponibles} en un único diccionario.

\begin{lstlisting}
dataset = loader.cargar_todo()
# Keys: metadata, fatiga, markers, chest, wrist, eeg, ear, nback, crt, task_switch
\end{lstlisting}

\subsection{Estructura de archivos soportados}

El módulo reconoce 43 archivos distribuidos en 6 categorías de sensores:

\begin{itemize}
    \item \textbf{Chest} (7 archivos): ECG, Acc, Respiración, RR intervals
    \item \textbf{Wrist/Empatica E4} (6 archivos): HR, EDA, BVP, Temperatura, Acc, IBI
    \item \textbf{Forehead/Muse} (8 archivos): 5 bandas EEG + Raw + Acc + Gyro
    \item \textbf{Ear/Nokia eSense} (6 archivos): PPG + Acc + Gyro (bilateral)
    \item \textbf{Zephyr} (2 archivos): Activity summary, Device status
    \item \textbf{Experimental} (5 archivos): Fatiga (VAS), Markers, N-Back, CRT, Task-Switch
\end{itemize}

% ========================================================================
\section{Módulo \texttt{processor.py}: Procesamiento de Series Temporales}
% ========================================================================

\subsection{Propósito}

Procesa datos cargados por \texttt{FatigueSetLoader} para extraer agregaciones y métricas por fase.

\subsection{Clase \texttt{DataProcessor}}

Procesador genérico de datos tabulares.

\begin{itemize}
    \item \texttt{process\_data(data)}: Suma todas las columnas numéricas.
    \item \texttt{filter\_data(data, threshold)}: Filtra filas por umbral.
    \item \texttt{aggregate\_data(data)}: Convierte dict a DataFrame.
\end{itemize}

\subsection{Clase \texttt{FatigueSetProcessor}}

Procesador especializado para FatigueSet.

\subsubsection{Inicialización}

\begin{lstlisting}
processor = FatigueSetProcessor(
    intensidad_map=loader.construir_intensidad_map()
)
\end{lstlisting}

\subsubsection{Método \texttt{calcular\_fatigabilidad()}}

Calcula deltas (cambios) de fatiga entre fases.

\begin{lstlisting}
df_fatiga_deltas = processor.calcular_fatigabilidad(df_fatiga)
# Output: 36 filas (12 participantes × 3 sesiones)
# Columnas: delta_fisica_ejercicio, delta_mental_ejercicio,
#          delta_fisica_cognitivo, delta_mental_cognitivo, etc.
\end{lstlisting}

La salida incluye:
\begin{itemize}
    \item \textbf{Deltas}: M2-M1, M3-M2, M3-M1 (para física y mental)
    \item \textbf{Scores absolutos}: Valores de fatiga en cada fase M1, M2, M3
    \item \textbf{Metadatos}: Participante, sesión, intensidad
\end{itemize}

\subsubsection{Método \texttt{segmentar\_por\_fase()}}

Divide una serie temporal en 3 tercios y retorna el segmento especificado.

\begin{lstlisting}
df_m1 = processor.segmentar_por_fase(df_todo, fase_num=0)  # Primeras 1/3 filas
df_m2 = processor.segmentar_por_fase(df_todo, fase_num=1)  # Segundas 1/3 filas
df_m3 = processor.segmentar_por_fase(df_todo, fase_num=2)  # Últimas 1/3 filas
\end{lstlisting}

\subsubsection{Métodos de Extracción de Métricas}

\textbf{Señales fisiológicas}:
\begin{lstlisting}
metricas = processor.metricas_segmento(df, ['hr', 'br'], prefijo='chest_')
# Output: {'chest_hr_media': ..., 'chest_hr_std': ..., ...}
\end{lstlisting}

\textbf{EEG por banda}:
\begin{lstlisting}
metricas_alpha = processor.metricas_eeg_segmento(df_eeg, banda='alpha')
# Output: {'eeg_alpha_media': valor_medio_4_canales}
\end{lstlisting}

\textbf{Tareas cognitivas}:
\begin{lstlisting}
metricas_nback = processor.metricas_nback(df_nback_fase)
# Output: {'nback_accuracy': %, 'nback_rt': ms}

metricas_crt = processor.metricas_crt(df_crt_fase)
# Output: {'crt_rt': ms, 'crt_accuracy': %}
\end{lstlisting}

\subsubsection{Método \texttt{construir\_dataset\_ml()}}

\textbf{Método central} que agrega todos los datos en un dataset de 108 filas.

\begin{lstlisting}
df_ml = processor.construir_dataset_ml(
    df_fatiga=dataset['fatiga'],
    df_chest=dataset['chest'],
    df_wrist_eda=dataset['wrist']['eda'],
    eeg_data=dataset['eeg'],
    df_nback=dataset['nback'],
    df_crt=dataset['crt'],
)
# Output: 108 filas (12 × 3 sesiones × 3 fases)
# Columnas: metadatos + features por modalidad
\end{lstlisting}

La salida estructura cada fila con:
\begin{itemize}
    \item Identificadores: \texttt{participante}, \texttt{sesion}, \texttt{fase}, \texttt{intensidad}
    \item Fatiga subjetiva: \texttt{fatiga\_fisica}, \texttt{fatiga\_mental}
    \item Features por sensor: media y std de cada modalidad
    \item Tareas cognitivas: accuracy y RT
\end{itemize}

% ========================================================================
\section{Módulo \texttt{features.py}: Extracción de Características}
% ========================================================================

\subsection{Propósito}

Extrae características estadísticas y especializadas de series temporales multimodales.

\subsection{Clase \texttt{FeatureExtractor}}

\subsubsection{Inicialización}

\begin{lstlisting}
extractor = FeatureExtractor(percentiles=[25, 50, 75])
\end{lstlisting}

\subsubsection{Métodos de Estadísticas Generales}

\begin{itemize}
    \item \texttt{extract\_features(data)}: Media de cada lista en un diccionario.
    \item \texttt{compute\_statistics(data)}: Estadísticas (mean, std, min, max) por clave.
\end{itemize}

\subsubsection{Método \texttt{features\_serie()}}

Estadísticas descriptivas de una serie temporal.

\begin{lstlisting}
feats = extractor.features_serie(serie, prefijo='hr_')
# Output: {
#    'hr_media': media,
#    'hr_std': desviación estándar,
#    'hr_min': mínimo,
#    'hr_max': máximo,
#    'hr_rango': max - min,
#    'hr_cv': coeficiente de variación,
#    'hr_p25': percentil 25, 'hr_p50': mediana, 'hr_p75': percentil 75
# }
\end{lstlisting}

\subsubsection{Método \texttt{features\_eeg\_banda()}}

Extrae características EEG por banda y canal.

\begin{lstlisting}
feats = extractor.features_eeg_banda(df_eeg_alpha, banda='alpha')
# Output: features por cada canal (TP9, AF7, AF8, TP10) + promedio global
\end{lstlisting}

\subsubsection{Método \texttt{features\_hrv()}}

Heart Rate Variability (HRV) a partir de intervalos RR.

\begin{lstlisting}
feats = extractor.features_hrv(df_rr_intervals)
# Output: {
#    'hrv_rmssd': raíz cuadrada del promedio de diferencias RR²,
#    'hrv_sdnn': desviación estándar de RR,
#    'hrv_pnn50': % de intervalos RR diferidos > 50ms,
#    'hrv_mean_rr': media de intervalos RR
# }
\end{lstlisting}

\subsubsection{Método \texttt{features\_eda()}}

Electrodermal Activity (EDA) - análisis de respuesta dérmica.

\begin{lstlisting}
feats = extractor.features_eda(df_eda)
# Output: {
#    'eda_media', 'eda_std', 'eda_min', 'eda_max', 'eda_cv': Estadísticas,
#    'eda_pendiente': tendencia lineal (cambio en tiempo),
#    'eda_n_picos_scr': número de picos de respuesta (SCR peaks)
# }
\end{lstlisting}

Los picos se detectan como cruces por encima de $\mu + \sigma$.

\subsubsection{Método \texttt{features\_cognitivas()}}

Características de tareas cognitivas (N-Back y CRT).

\begin{lstlisting}
feats = extractor.features_cognitivas(
    df_nback=df_nback_fase,
    df_crt=df_crt_fase
)
# Output: {
#    'nback_accuracy': %, 'nback_rt_media': ms, 'nback_rt_std': ms,
#    'nback_n_errores': número de errores,
#    'crt_rt_media': ms, 'crt_rt_std': ms
# }
\end{lstlisting}

% ========================================================================
\section{Módulo \texttt{validators.py}: Validación de Integridad}
% ========================================================================

\subsection{Propósito}

Valida la integridad del dataset FatigueSet (valores ausentes, duplicados, tipos mixtos).

\subsection{Clase \texttt{DataValidator}}

Validador genérico de DataFrames.

\begin{itemize}
    \item \texttt{validate\_schema(df, columns)}: Comprueba presencia de todas las columnas.
    \item \texttt{check\_missing\_values(df)}: Detecta valores NaN.
    \item \texttt{validate\_data\_quality(df)}: Verifica coherencia de tipos de datos.
\end{itemize}

\subsection{Clase \texttt{FatigueSetValidator}}

Validador especializado para FatigueSet.

\subsubsection{Inicialización}

\begin{lstlisting}
validator = FatigueSetValidator(
    loader=loader,
    umbral_nulos=5.0  # % máximo permitido
)
\end{lstlisting}

\subsubsection{Método \texttt{validar\_todos()}}

Valida todos los archivos críticos (11 tipos de sensores + 3 tareas).

\begin{lstlisting}
df_validacion = validator.validar_todos()
# Output: DataFrame con una fila por (archivo, participante, sesión)
# Columnas: Archivo, Participante, Sesión, Filas, Nulos (%), Duplicados, Estado
\end{lstlisting}

Archivos críticos validados:
\begin{enumerate}
    \item Chest Physiology, ECG Raw
    \item Wrist HR, EDA, Temp
    \item EEG Alpha, Beta
    \item Ear PPG Left
    \item Fatigue Scores, N-Back Task, CRT Task
\end{enumerate}

\subsubsection{Método \texttt{resumen()}}

Genera agregaciones de problemas por tipo de archivo.

\begin{lstlisting}
resumen_df, problemas_df = validator.resumen(df_validacion)
# resumen_df: resumen por archivo (sesiones OK, nulos máximos)
# problemas_df: sesiones con estado '⚠️'
\end{lstlisting}

\subsubsection{Método \texttt{imprimir\_resumen()}}

Imprime un reporte legible en consola.

% ========================================================================
\section{Módulo \texttt{utils.py}: Funciones Auxiliares}
% ========================================================================

\subsection{Función \texttt{convertir\_timestamp()}}

Convierte columnas de timestamp (milisegundos desde epoch) a \texttt{datetime} de pandas.

\begin{lstlisting}
df = convertir_timestamp(df, col='timestamp')
# Crea columna 'datetime' con formato datetime64
\end{lstlisting}

% ========================================================================
\section{Flujo de Trabajo Completo}
% ========================================================================

\subsection{Pipeline End-to-End}

Un uso típico completo:

\begin{lstlisting}
from fatigueset.loader import FatigueSetLoader
from fatigueset.processor import FatigueSetProcessor
from fatigueset.features import FeatureExtractor
from fatigueset.validators import FatigueSetValidator

# 1. CARGA
loader = FatigueSetLoader(base_path='fatigueset')
dataset = loader.cargar_todo()

# 2. VALIDACIÓN
validator = FatigueSetValidator(loader)
validator.imprimir_resumen()

# 3. PROCESAMIENTO
processor = FatigueSetProcessor(
    intensidad_map=loader.construir_intensidad_map()
)
df_ml = processor.construir_dataset_ml(
    df_fatiga=dataset['fatiga'],
    df_chest=dataset['chest'],
    df_wrist_eda=dataset['wrist']['eda'],
    eeg_data=dataset['eeg'],
    df_nback=dataset['nback'],
    df_crt=dataset['crt'],
)

# 4. EXTRACCIÓN DE FEATURES AVANZADAS
extractor = FeatureExtractor()
for banda in ['alpha', 'beta', 'theta']:
    df_ml[f'eeg_{banda}_advanced'] = df_ml.apply(
        lambda row: extractor.features_eeg_banda(
            dataset['eeg'][banda], banda
        ), axis=1
    )

# Dataset final: 108 × ~50 columnas
print(df_ml.shape)  # (108, 50+)
df_ml.to_csv('dataset_fatigueset_ml.csv', index=False)
\end{lstlisting}

\subsection{Etapas del pipeline}

\textbf{Etapa 1: Carga}
\begin{itemize}
    \item Lectura de 43 archivos CSV distribuidos
    \item Sincronización de metadatos (intensidades)
    \item Consolidación en estructura de diccionarios
\end{itemize}

\textbf{Etapa 2: Validación}
\begin{itemize}
    \item Comprobación de 11 tipos de sensores
    \item Detección de valores nulos y duplicados
    \item Reporte de problemas por sesión
\end{itemize}

\textbf{Etapa 3: Procesamiento}
\begin{itemize}
    \item Segmentación temporal en 3 fases por sesión
    \item Extracción de métricas sumarias (media, std)
    \item Cálculo de deltas de fatiga (M2-M1, M3-M2, etc.)
\end{itemize}

\textbf{Etapa 4: Features}
\begin{itemize}
    \item Features especializadas por modalidad (HRV, EDA, EEG)
    \item Características de tareas cognitivas (accuracy, RT)
    \item Preparación de variables para ML
\end{itemize}

\textbf{Etapa 5: Aggregación}
\begin{itemize}
    \item Dataset final 108 × 50+ columnas
    \item Listo para modelos predictivos
\end{itemize}

% ========================================================================
\section{Integración con Procesamiento de Series Temporales}
% ========================================================================

Esta librería implementa en Python los conceptos descritos en el documento \textit{``Procesamiento de Datos de Series Temporales''}.

\subsection{Sincronización Multi-Sensor}

Problema: Sensores con diferentes frecuencias de muestreo y relojes internos.

\begin{center}
\begin{tabular}{|l|c|}
\toprule
\textbf{Sensor} & \textbf{Frecuencia (Hz)} \\
\midrule
Chest HR & 1 \\
Wrist EDA & 4 \\
Forehead EEG & 9.13 \\
\bottomrule
\end{tabular}
\end{center}

Solución: \texttt{loader.cargar\_todo()} sincroniza mediante timestamps en \texttt{exp\_markers.csv}.

\subsection{Windowing y Segmentación}

En lugar de ventanas deslizantes explícitas, la librería segmenta cada sesión en 3 phases por tercio temporal:

$$\text{fila} = \lfloor \frac{\text{fase\_num} \times n}{n\_fases} \rfloor, \quad \text{fila} = \lfloor \frac{(\text{fase\_num}+1) \times n}{n\_fases} \rfloor$$

Método: \texttt{FatigueSetProcessor.segmentar\_por\_fase()}

\subsection{Data Leakage Prevention}

La librería implementa normalización sin contaminación:
\begin{enumerate}
    \item Separación temporal: Split determinístico por fases
    \item Normalización independiente por segmento si es necesario
    \item Metadatos claros (participante, sesión, fase)
\end{enumerate}

% ========================================================================
\section{Estructura del Dataset ML Agregado}
% ========================================================================

\subsection{Dimensiones}

\begin{itemize}
    \item \textbf{Filas}: 108 (12 participantes $\times$ 3 sesiones $\times$ 3 fases)
    \item \textbf{Columnas}: Aproximadamente 50+, distribuidas como:
        \begin{itemize}
            \item 6 columnas de metadatos
            \item 15-20 features fisiológicas
            \item 10-15 features EEG
            \item 8-10 features cognitivas
            \item 5-10 features especializadas (HRV, EDA)
        \end{itemize}
\end{itemize}

\subsection{Columnas principales}

\subsubsection{Metadatos}
\begin{lstlisting}
participante, sesion, fase, intensidad, intensidad_num, fase_num
\end{lstlisting}

\subsubsection{Fatiga Subjetiva (VAS)}
\begin{lstlisting}
fatiga_fisica, fatiga_mental
\end{lstlisting}

\subsubsection{Señales Fisiológicas}
\begin{lstlisting}
chest_hr_media, chest_hr_std, chest_br_media, chest_br_std, chest_hrv_media
eda_media, eda_std, eda_pendiente, eda_n_picos_scr
\end{lstlisting}

\subsubsection{EEG}
\begin{lstlisting}
eeg_alpha_media, eeg_beta_media, eeg_theta_media
eeg_alpha_TP9_media, eeg_alpha_AF7_media, ... (por canal)
\end{lstlisting}

\subsubsection{Tareas Cognitivas}
\begin{lstlisting}
nback_accuracy, nback_rt_media, nback_rt_std, nback_n_errores
crt_rt_media, crt_rt_std
\end{lstlisting}

% ========================================================================
\section{Anexos}
% ========================================================================

\subsection{Tabla de Sensores y Frecuencias}

\begin{center}
\begin{tabular}{|l|l|c|}
\toprule
\textbf{Dispositivo} & \textbf{Sensor} & \textbf{Frecuencia (Hz)} \\
\midrule
\multirow{3}{*}{Zephyr BioHarness} & Heart Rate & 1 \\
& Breathing Rate & 1 \\
& Corrected GSR & 1 \\
\midrule
\multirow{6}{*}{Empatica E4} & HR (IBI) & 1 \\
& EDA & 4 \\
& BVP & 64 \\
& Temperatura piel & 4 \\
& Acelerómetro & 32 \\
\midrule
\multirow{5}{*}{InteraXon Muse} & EEG TP9, AF7, AF8, TP10 & 9.13 \\
& 5 bandas (alpha, beta, delta, gamma, theta) & 1 \\
& PPG (Ear channels) & 64 \\
\bottomrule
\end{tabular}
\end{center}

\subsection{Estructura de Carpetas del Dataset}

\begin{lstlisting}
fatigueset/
├── 01/
│   ├── 01/  (sesión 1)
│   │   ├── chest_*.csv
│   │   ├── wrist_*.csv
│   │   ├── forehead_eeg_*.csv
│   │   ├── exp_fatigue.csv
│   │   ├── exp_nback.csv
│   │   └── exp_crt.csv
│   ├── 02/  (sesión 2)
│   └── 03/  (sesión 3)
├── 02/ ... 12/
├── metadata.csv  (mapeo participante → intensidades)
├── pre_task_survey.xlsx
└── preliminary_questionnaire.xlsx
\end{lstlisting}

\subsection{Referencia de Constantes}

\begin{lstlisting}
# Mapa de fases
FASES_MEDICION = {
    0: 'M1_baseline',
    1: 'M2_post_ejercicio',
    2: 'M3_post_fatiga_mental'
}

# Mapa de intensidades
INTENSIDAD_NUM = {'low': 1, 'medium': 2, 'high': 3}

# Canales EEG del Muse
CANALES_EEG = ['TP9', 'AF7', 'AF8', 'TP10']
\end{lstlisting}

\subsection{Notebooks de Ejemplo}

La librería incluye ejemplos prácticos:

\begin{itemize}
    \item \texttt{notebooks/example\_usage.ipynb}: Demostración completa del pipeline
    \item \texttt{../Jupyters/FatigueSet Procesado de datos.ipynb}: Pipeline original del proyecto
    \item \texttt{../Jupyters/windowing\_data\_leakage.ipynb}: Análisis de data leakage
    \item \texttt{../Jupyters/sincronizacion\_normalizacion.ipynb}: Sincronización multi-sensor
\end{itemize}

\subsection{Referencias Cruzadas}

Para más detalles sobre conceptos de procesamiento de series temporales:

\begin{itemize}
    \item \texttt{../series-temporales/procesamiento\_datos\_series\_temporales.pdf}
    \item Secciones: Sincronización, Windowing, Data Leakage, Normalización
\end{itemize}

% ========================================================================
\appendix

\section{Manejo de Errores Comunes}

\subsection{Error: ``No se encontró el dataset''}

\begin{lstlisting}
FileNotFoundError: No se encontró el dataset en: fatigueset
\end{lstlisting}

\textbf{Solución}: Asegurar que la ruta \texttt{base\_path} es correcta y que la carpeta contiene las subcarpetas 01--12.

\subsection{Error: ``Metadata.csv no encontrado''}

\textbf{Solución}: Verificar que existe \texttt{fatigueset/metadata.csv}.

\subsection{Advertencia: ``Archivo no encontrado en P01/S02''}

Esto es esperado si no todos los archivos están disponibles para todas las sesiones. La librería continúa procesando.

\end{document}
